% !TITLE 破阵子·四十年来家国
% !AUTHOR 李煜
% !DYNASTY 五代十国
% !GENRE 词
% !ORDER 3

\begin{poem}
四十年来家国\textsuperscript{1}，三千里地山河\textsuperscript{2}。
凤阁龙楼\textsuperscript{3}连霄汉\textsuperscript{4}，玉树琼枝\textsuperscript{5}作烟萝\textsuperscript{6}，几曾识干戈\textsuperscript{7}？

一旦归为臣虏\textsuperscript{8}，沈腰潘鬓\textsuperscript{9}消磨。
最是仓皇辞庙\textsuperscript{10}日，教坊\textsuperscript{11}犹奏别离歌，垂泪对宫娥\textsuperscript{12}。
\end{poem}

\begin{annotations}
\notetext{1}{四十年家国：南唐自937年立国到975年灭亡，前后约四十年。开篇以宏大的数字起笔，写出故国的久远与辽阔。}
\notetext{2}{三千里地山河：南唐疆域包括今天的江苏、安徽、江西等地，方圆约三千里。先写时间，再写空间，气象雄伟。}
\notetext{3}{凤阁龙楼：雕饰着龙凤的楼阁，指皇宫建筑。凤和龙是帝王的象征。}
\notetext{4}{霄汉：云霄和天河，指极高的天空。皇宫楼阁高耸入云，可见昔日南唐宫殿的壮丽。}
\notetext{5}{玉树琼枝：像玉一样的树木，像美玉一样的枝条。琼（qióng），红色的美玉。形容宫中花木名贵华美。}
\notetext{6}{烟萝：烟雾中的藤萝。把宫中花木看作山野烟萝——说明李煜在宫中过着与世隔绝的生活，不知人间疾苦。}
\notetext{7}{干戈：盾牌和长戈，代指战争。从小生长在深宫中，什么时候见过战争呢？这句自问，是对自己前半生的反思。}
\notetext{8}{臣虏：被俘虏的臣子。从一国之君沦为阶下囚，“一旦”二字写尽命运的突然。}
\notetext{9}{沈腰潘鬓：沈约（南朝文学家）曾说自己腰带一天天变松（消瘦）；潘岳（西晋诗人）在作品中写自己三十二岁就鬓发斑白。后世用“沈腰”指消瘦，“潘鬓”指白发。李煜用这个典故说自己日渐憔悴衰老。}
\notetext{10}{辞庙：告别宗庙。宗庙是帝王祭祀祖先的地方，辞庙意味着永远告别自己的国家。}
\notetext{11}{教坊：朝廷掌管音乐的机构。临别时教坊还在演奏别离的曲子——音乐没有错，但听在耳里，全是撕心裂肺的痛。}
\notetext{12}{宫娥：宫女。亡国的最后时刻，他流着泪对着宫女，而不是对着江山社稷。这被认为是李煜“不争气”的表现——但正是这种儿女情长，让他的词格外动人。}
\end{annotations}

\begin{appreciation}
写亡国时刻的词，这是最有名的一首之一。亡国这个词，我们平时只在历史书里见到，这一首把它变成了一滴看得见的眼泪。

上片全是数字和气势：“四十年来家国，三千里地山河”——四十年，是时间；三千里，是空间。两个大数字，把南唐的家底一下子铺开。“凤阁龙楼连霄汉”，宫殿高耸入云；“玉树琼枝作烟萝”，宫里的奇花异木，被他看得像山野间的藤萝一样寻常。然后是一个反问：“几曾识干戈？”什么时候见过战争？一辈子没摸过兵器，所以亡国亡得毫无防备。

下片“一旦”两个字，命运的转折快得吓人：“一旦归为臣虏，沈腰潘鬓消磨。”沈腰是南朝沈约，腰一天天变细；潘鬓是西晋潘岳，年纪轻轻白了头。两个典故合起来就一句话：人瘦了，也老了。被俘之后的日子，是把人一天天磨掉的。

“最是仓皇辞庙日，教坊犹奏别离歌，垂泪对宫娥。”辞庙，是去宗庙告别祖先的牌位，从此永远离开自己的国家。临别时，宫里的乐队还在奏离别的曲子——音乐没有错，听在耳里却全是撕心裂肺的痛。最后他对着宫女哭，而不是对着江山社稷。有人说他不争气，可正是这份儿女情长，让他不像个帝王，而像个人。屈原在《离骚》里哀的是“民生之多艰”，李煜哀的是自己；一个为国，一个为己，但两个人都把自己的心整个放进作品里，这是中国抒情诗的正路。《采薇》里那位战士说“我心伤悲，莫知我哀”——李煜的悲哀，同样没人真正懂，他只能垂泪对宫娥。

顺便一提，破阵子这个词牌，后来辛弃疾也填过一首壮词，写的是梦里杀敌——同一个曲调，一个填回不去的战场，一个填回不去的家国。

李煜这首词写完，再也没有写过“三千里地山河”这么大的词，他后来的词里只剩下小楼、西楼和一江春水。人失去得太多，话就少了。你有心事的时候别学他闷着，想说就说，哥哥听着。
\end{appreciation}
